\documentclass[11pt]{article}
% ===========================================================================
%  Shared preamble for the AP Calculus AB exam-preparation worksheets.
%  Input right after \documentclass. Files beginning with "_" are not compiled
%  on their own. Navy and gold, to match the course's lesson pages.
% ===========================================================================
\usepackage[letterpaper,top=2.3cm,bottom=2.2cm,left=1.7cm,right=1.7cm,headheight=28pt,footskip=30pt]{geometry}
\usepackage{amsmath,amssymb}
\usepackage[T1]{fontenc}
\usepackage[utf8]{inputenc}
\usepackage{lmodern}
\usepackage{xcolor}
\usepackage[most]{tcolorbox}
\usepackage{enumitem}
\usepackage{fancyhdr}
\usepackage{array}
\usepackage{tikz}
\usepackage{pgfplots}
\pgfplotsset{compat=1.18}

\definecolor{navy}{HTML}{161E54}
\definecolor{ink}{HTML}{27357E}
\definecolor{gold}{HTML}{B98A00}
\definecolor{goldbg}{HTML}{FFFDF3}
\definecolor{inkbg}{HTML}{F6F7FD}
\definecolor{rule}{HTML}{C9CFF0}
\definecolor{okgreen}{HTML}{1B7A44}
\definecolor{okbg}{HTML}{F1F8F3}

\pagestyle{fancy}
\fancyhf{}
\fancyhead[L]{\footnotesize Name: \underline{\hspace{4cm}}}
\fancyhead[C]{\textbf{\color{navy}Integration Academy}\\[-2pt]{\scriptsize Calculus AB \textperiodcentered{} AP\textsuperscript{\textregistered} Exam Preparation}}
\fancyfoot[L]{\tiny\color{black!55}\parbox[t]{0.78\textwidth}{AP\textsuperscript{\textregistered} is a registered trademark of the College Board, which was not involved in the production of, and does not endorse, this material. \textcopyright{} Integration Academy.}}
\fancyfoot[R]{\scriptsize\thepage}
\renewcommand{\headrulewidth}{1.2pt}
\renewcommand{\headrule}{\hbox to\headwidth{\color{gold}\leaders\hrule height \headrulewidth\hfill}}
\renewcommand{\footrulewidth}{0pt}
\setlength{\parindent}{0pt}
\setlength{\parskip}{4pt}

% Title block: \wstitle{code}{title}{unit line}
\newcommand{\wstitle}[3]{\begin{center}{\small\bfseries\color{gold}WORKSHEET #1}\\[3pt]{\LARGE\bfseries\color{navy}#2}\\[3pt]{\small\color{black!60}#3}\end{center}\vspace{2pt}}
% Section heading
\newcommand{\sech}[1]{\par\vspace{10pt}{\large\bfseries\color{navy}#1}\par\nobreak\vspace{2pt}{\color{rule}\hrule height0.8pt}\vspace{7pt}}
% A level tag printed beside a group of questions
\newcommand{\level}[1]{\par\vspace{6pt}{\footnotesize\bfseries\color{gold}\MakeUppercase{#1}}\par\nobreak\vspace{2pt}}

\newtcolorbox{ideabox}{colframe=ink,colback=inkbg,coltitle=white,colbacktitle=ink,fonttitle=\bfseries\small,title={Key ideas},arc=2pt,boxrule=0.8pt,left=8pt,right=8pt}
\newtcolorbox{example}[1]{colframe=rule,colback=white,coltitle=navy,colbacktitle=inkbg,fonttitle=\bfseries\small,title={#1},arc=2pt,boxrule=0.8pt,borderline west={3pt}{0pt}{ink},left=10pt,right=8pt}
\newtcolorbox{frbox}[1]{colframe=gold,colback=goldbg,coltitle=white,colbacktitle=gold,fonttitle=\bfseries\small,title={#1},arc=2pt,boxrule=0.8pt,breakable,left=8pt,right=8pt}
\newtcolorbox{solbox}{colframe=okgreen,colback=okbg,coltitle=white,colbacktitle=okgreen,fonttitle=\bfseries\small,title={Worked solutions},arc=2pt,boxrule=0.8pt,breakable,left=8pt,right=8pt}
\newtcolorbox{scorebox}[1]{colframe=navy,colback=white,coltitle=white,colbacktitle=navy,fonttitle=\bfseries\small,title={#1},arc=2pt,boxrule=0.8pt,breakable,left=8pt,right=8pt}

\newcommand{\soln}{\par\smallskip\textbf{\color{ink}Solution.}\ }
\newcommand{\result}[1]{\par\smallskip\textbf{#1}}
% blank working space of a given height
\newcommand{\work}[1]{\par\vspace{#1}}
% a point value in the margin of a scoring guide
\newcommand{\pts}[1]{\hfill{\footnotesize\bfseries\color{gold}[#1]}}

% Coordinate grid for graph questions: \begin{apgraph}{xmin}{xmax}{ymin}{ymax} ... \end{apgraph}
\newenvironment{apgraph}[4]{\begin{center}\begin{tikzpicture}\begin{axis}[
  axis lines=middle, xmin=#1, xmax=#2, ymin=#3, ymax=#4,
  grid=both, grid style={black!12}, xtick distance=1, ytick distance=1,
  ticklabel style={font=\scriptsize}, xlabel={$x$}, ylabel={$y$},
  width=10cm, height=7.2cm, samples=80, clip=false,
  every axis plot/.append style={line width=1.1pt, color=ink}]}{\end{axis}\end{tikzpicture}\end{center}}
\newcommand{\closeddot}[2]{\addplot[only marks, mark=*, mark size=2.2pt, color=ink] coordinates {(#1,#2)};}
\newcommand{\opendot}[2]{\addplot[only marks, mark=*, mark size=2.2pt, color=ink, mark options={fill=white}] coordinates {(#1,#2)};}

% Multiple-choice options: four across, or stacked when the options are long.
\newcommand{\opts}[4]{\par\vspace{3pt}\begin{tabular}{@{}p{0.235\linewidth}p{0.235\linewidth}p{0.235\linewidth}p{0.235\linewidth}@{}}(A)\ #1 & (B)\ #2 & (C)\ #3 & (D)\ #4\end{tabular}\par\vspace{9pt}}
\newcommand{\optsv}[4]{\par\vspace{3pt}\begin{tabular}{@{}l@{\ \ }p{0.9\linewidth}@{}}(A) & #1\\(B) & #2\\(C) & #3\\(D) & #4\end{tabular}\par\vspace{9pt}}

\fancyhead[R]{\footnotesize\color{black!60}Worksheet 7.5}
\begin{document}

\wstitle{7.5}{Separation of Variables}{Unit 7: Differential Equations}

\begin{ideabox}
\begin{itemize}[leftmargin=*,itemsep=2pt,topsep=2pt]
  \item If $\dfrac{dy}{dx}=g(x)\,h(y)$, move everything in $y$ to one side and everything in $x$ to the other: $\dfrac{dy}{h(y)}=g(x)\,dx$.
  \item Integrate both sides. Add \textbf{one} constant of integration, at this step.
  \item Then solve for $y$. An exponential turns the constant into a multiplier: $y=Ae^{\ldots}$.
  \item Dividing by $h(y)$ can lose a constant solution where $h(y)=0$. Check for it separately.
\end{itemize}
\end{ideabox}

\sech{Worked examples}
\begin{example}{Example 1 \textperiodcentered{} An exponential solution}
Solve $\dfrac{dy}{dx}=3x^2y$.\soln $\dfrac{dy}{y}=3x^2\,dx$, so $\ln|y|=x^3+C$.\result{$y=Ae^{x^3}$.}
\end{example}

\begin{example}{Example 2 \textperiodcentered{} A square root}
Solve $\dfrac{dy}{dx}=\dfrac{\cos x}{y}$.\soln $y\,dy=\cos x\,dx$, so $\dfrac{y^2}{2}=\sin x+C_0$.\result{$y^2=2\sin x+C$.}
\end{example}

\begin{example}{Example 3 \textperiodcentered{} A reciprocal}
Solve $\dfrac{dy}{dx}=\dfrac{y^2}{x}$ for $x>0$, $y\neq0$.\soln $y^{-2}\,dy=\dfrac{dx}{x}$, so $-\dfrac1y=\ln x+C$.\result{$y=-\dfrac{1}{\ln x+C}$.}
\end{example}

\newpage
\sech{Practice}

\level{Core}
\begin{enumerate}[leftmargin=*,itemsep=2pt]
  \item Separate the variables in $\dfrac{dy}{dx}=xy^3$. Do not integrate.\work{2.3cm}
  \item Solve $\dfrac{dy}{dx}=4y$.\work{2.3cm}
  \item Solve $\dfrac{dy}{dx}=6x^2$.\work{2.3cm}
  \item Solve $\dfrac{dy}{dx}=-y$.\work{2.3cm}
\end{enumerate}

\level{Exam level}
\begin{enumerate}[leftmargin=*,itemsep=2pt,start=5]
  \item Solve $\dfrac{dy}{dx}=\dfrac{x+1}{y}$, leaving your answer in the form $y^2=\ldots$\work{2.8cm}
  \item Solve $\dfrac{dy}{dx}=y\sin x$.\work{2.8cm}
  \item Solve $\dfrac{dy}{dx}=e^{2x-y}$.\work{2.8cm}
  \item Solve $\dfrac{dy}{dx}=\sqrt{1-y^2}$ for $-1<y<1$.\work{2.8cm}
\end{enumerate}

\begin{frbox}{Question 9 \textperiodcentered{} Free-response practice}
Consider the differential equation $\dfrac{dy}{dx}=x\sqrt{y}$ for $y\ge0$.
\begin{enumerate}[label=(\alph*),leftmargin=*,itemsep=2pt]
  \item Separate the variables.\work{2.6cm}
  \item Find the general solution for $y>0$, in the form $y=\ldots$\work{2.6cm}
  \item Show that $y=0$ is also a solution, and explain why the method of part (b) did not produce it.\work{2.6cm}
\end{enumerate}
\end{frbox}

\newpage
\begin{solbox}
\begin{enumerate}[leftmargin=*,itemsep=4pt]
  \item $\dfrac{dy}{y^3}=x\,dx$.
  \item $y=Ae^{4x}$.
  \item $y=2x^3+C$.
  \item $y=Ae^{-x}$.
  \item $y\,dy=(x+1)\,dx$: $y^2=x^2+2x+C$.
  \item $\ln|y|=-\cos x+C$: $y=Ae^{-\cos x}$.
  \item $e^{y}\,dy=e^{2x}\,dx$: $e^{y}=\dfrac12e^{2x}+C$, so $y=\ln\!\left(\dfrac12e^{2x}+C\right)$.
  \item $\dfrac{dy}{\sqrt{1-y^2}}=dx$: $\arcsin y=x+C$, so $y=\sin(x+C)$.
  \item \begin{enumerate}[label=(\alph*),leftmargin=*,itemsep=3pt]
    \item $\dfrac{dy}{\sqrt y}=x\,dx$.
    \item $2\sqrt y=\dfrac{x^2}{2}+C_0$, so $y=\left(\dfrac{x^2}{4}+C\right)^2$.
    \item If $y=0$ for all $x$, both sides of the equation are 0. Separating required division by $\sqrt y$, which is not allowed when $y=0$.
  \end{enumerate}
\end{enumerate}
\end{solbox}

\end{document}
